\section{Technical design}
\label{sec:design}

This section states the algorithm as \ttlibname implements it. The notation throughout:
$B$ examples per batch, $F$ Boolean features, $2F$ literals, $C$ clauses, $K$ outputs,
$N$ automaton states per action, $s$ specificity, $T$ vote margin.

\subsection{Representation}

An example is a Boolean vector $x \in \{0,1\}^F$. Its \emph{literal vector} is
$L(x) = [x, \lnot x] \in \{0,1\}^{2F}$. Every (clause, literal) pair owns a Tsetlin automaton
whose state $\theta_{jk} \in \{0,\dots,2N-1\}$ decides membership:

\[
  \mathrm{include}_{jk} \;=\; \mathbb{1}\!\left[\theta_{jk} \ge N\right],
  \qquad
  c_j(x) \;=\; \bigwedge_{k \,:\, \mathrm{include}_{jk}} L_k(x).
\]

The whole learned state is therefore one integer matrix $\theta \in \mathbb{Z}^{C \times 2F}$
--- in \ttlibname, the \code{int32} buffer \code{ta\_state}. Two derived, non-persistent
buffers cache the include mask as floats and the per-clause literal count; anything that
mutates $\theta$ refreshes them, including \code{load\_state\_dict}.

\begin{figure}[htbp]
\centering
\resizebox{0.95\linewidth}{!}{% Figure: one Tsetlin automaton — the state line, the include boundary, and the transitions.
\begin{tikzpicture}[font=\sffamily\small]

  \def\cw{1.12}   % cell width

  % ---------------------------------------------------------------- the state line
  \foreach \i in {0,...,11}{
    \pgfmathsetmacro{\x}{\i*\cw}
    \ifnum\i<6
      \node[draw=ttgray!50, fill=white, minimum width=\cw cm, minimum height=0.8cm,
            inner sep=0pt, anchor=south west, font=\small] at (\x,0) {\i};
    \else
      \pgfmathtruncatemacro{\shade}{18+8*(\i-6)}
      \node[draw=ttdeep!60, fill=ttorange!\shade,
            minimum width=\cw cm, minimum height=0.8cm,
            inner sep=0pt, anchor=south west, font=\small] at (\x,0) {\i};
    \fi
  }

  % ---------------------------------------------------------------- boundary
  \draw[ttdeep, line width=1.6pt] (6*\cw,-0.18) -- (6*\cw,1.28);
  \node[tttitle, anchor=south, font=\sffamily\bfseries\small, text=ttdeep]
    at (6*\cw,1.32) {include boundary $N$};

  % ---------------------------------------------------------------- motion arrows
  \draw[ttarrow, line width=1.1pt] (1.6*\cw,1.62) -- (3.4*\cw,1.62);
  \node[ttlabel, anchor=south, text=ttorange, font=\sffamily\small]
    at (2.5*\cw,1.64) {memorise $+1$};
  \draw[ttarrow, line width=1.1pt, draw=ttslate] (10.4*\cw,1.62) -- (8.6*\cw,1.62);
  \node[ttlabel, anchor=south, text=ttslate, font=\sffamily\small]
    at (9.5*\cw,1.64) {forget $-1$};

  % ---------------------------------------------------------------- init marker
  \draw[ttarrow, draw=ttslate] (5.5*\cw,2.95) -- (5.5*\cw,0.88);
  \node[ttlabel, anchor=south, text=ttslate, align=center, font=\sffamily\small]
    at (5.5*\cw,2.98) {automata start at $N-1$ --- \emph{about to be memorised, currently forgotten}};

  % ---------------------------------------------------------------- half labels
  \node[ttlabel, anchor=north, text width=5.6cm, align=center, font=\sffamily\small] at (3*\cw,-0.28)
    {\textbf{excluded}: the literal is \emph{forgotten}, states $0 \dots N-1$};
  \node[ttlabel, anchor=north, text width=5.6cm, align=center, text=ttdeep, font=\sffamily\small] at (9*\cw,-0.28)
    {\textbf{included}: the literal is \emph{memorised}, states $N \dots 2N-1$};

  % ---------------------------------------------------------------- transitions
  \begin{scope}[shift={(0,-2.55)}]
    \node[ttnode, anchor=north west, text width=3.65cm] (ia) at (0,0)
      {\textbf{Type Ia} \textemdash\ \emph{Recognize}\\[2pt]
       \scriptsize the clause matched: True literals $+1$ with probability
       $\frac{s-1}{s}$ (or $1$ when boosted), False literals $-1$ with probability $\frac{1}{s}$};
    \node[ttnode, anchor=north west, text width=3.65cm] (ib) at (4.55,0)
      {\textbf{Type Ib} \textemdash\ \emph{Erase}\\[2pt]
       \scriptsize the clause did not match: every literal $-1$ with probability $\frac{1}{s}$};
    \node[ttnode, anchor=north west, text width=3.65cm] (ii) at (9.1,0)
      {\textbf{Type II} \textemdash\ \emph{Reject}\\[2pt]
       \scriptsize the clause matched an example of another class: excluded False literals
       $+1$, deterministically, never past $N$};
  \end{scope}

  \draw[ttarrow] (2.0,-1.55) -- (2.0,-2.5);
  \draw[ttarrow] (6.55,-1.55) -- (6.55,-2.5);
  \draw[ttarrow] (11.1,-1.55) -- (11.1,-2.5);
\end{tikzpicture}
}
\caption{One Tsetlin automaton. The state line is split at $N$: below it the literal is
  forgotten, at or above it the literal is included in the clause. Deep states are hard to
  dislodge, states near the boundary flip easily --- this is the memory that makes the
  machine robust to noisy examples.}
\label{fig:automaton}
\end{figure}

\subsection{Inference is two matrix products}
\label{sec:inference}

A conjunction is awkward to vectorise directly, but its negation is a sum. A clause fails on
an example exactly when it includes at least one literal that is False, so

\[
  \mathrm{violations} \;=\; (1 - L)\, \mathrm{include}^{\!\top} \in \mathbb{R}^{B\times C},
  \qquad
  c_j(x_b) \;=\; \mathbb{1}\!\left[\mathrm{violations}_{bj} = 0\right].
\]

One dense matrix product evaluates every clause on every example. The vote sums are a second
product, $v = c\,W$, clipped to $[-T, T]$, where $W \in \mathbb{Z}^{C\times K}$ holds the
signed clause weights. Prediction is $\hat y = \arg\max_k v_k$.

\begin{ttcallout}[Empty clauses are mode-dependent --- and this is not a detail]
A clause with no included literals has zero violations, so the formula above makes it match
\emph{everything}. That is correct while learning: a fresh clause must fire in order to
receive the Type~Ia feedback that gives it its first literal. It is wrong at prediction time,
where an empty clause should vote for nothing. \ttlibname keys this off \code{self.training},
exactly like \code{nn.Dropout}, so \code{model.eval()} is required before measuring accuracy
--- forgetting it silently changes predictions rather than raising. The behaviour is pinned
by tests.
\end{ttcallout}

\subsection{The feedback policy}

For a labelled example the machine decides, per clause, whether to give Type~I feedback
(make the clause recognise this example), Type~II feedback (make it stop matching), or
nothing. The decision is probabilistic in the vote margin: the target class receives
feedback with probability $(T - v_y)/2T$ and a sampled negative class with probability
$(T + v_q)/2T$. Once a class already wins by $T$ votes its clauses stop being updated, which
frees them to specialise elsewhere --- the resource-allocation mechanism at the heart of the
algorithm~\ttcite{1}.

Each model differs only in how votes are formed and who receives which feedback
(Table~\ref{tab:policies}). Everything below that line is shared.

\begin{table}[htbp]
\centering
\small\sffamily\sloppy
\caption{Output layer and feedback policy per model. This table \emph{is} the difference
  between the model classes; the learning machinery underneath is identical.}
\label{tab:policies}
\begin{tabularx}{\linewidth}{@{}L{3.5cm}L{4.15cm}>{\raggedright\arraybackslash}X@{}}
\ttoprule
\thead{Model} & \thead{Vote sums} & \thead{Who receives what}\\
\ttmidrule
\clsname{TsetlinMachine} &
  $v_k=\sum_{j\in\text{pos}(k)} w_jc_j-\sum_{j\in\text{neg}(k)} w_jc_j$ &
  target class: Type~I to positive, Type~II to negative clauses; one random negative class:
  the reverse\\
\clsname{CoalescedTsetlinMachine} &
  $v = c\,W$, $W\in\mathbb{Z}^{C\times K}$ &
  per output, Type~I when the weight sign agrees with the target and Type~II when it does
  not; weights move by $\pm1$, so a clause can change sides\\
\clsname{RegressionTsetlinMachine} &
  $v=\sum_j w_jc_j$ clipped to $[0,T]$, mapped to $[y_{\min},y_{\max}]$ &
  Type~I with probability $(v^\star-v)/T$ when the prediction is too low, Type~II with
  $(v-v^\star)/T$ when too high\\
\ttbotrule
\end{tabularx}
\end{table}

Two refinements are available on the classifiers. \emph{Integer clause weights} (the weighted
TM~\ttcite{3}) let one clause count for many, which is what makes large convolutional
machines tractable: Type~Ia increments a clause's weight, Type~II decrements it.
\emph{Focused negative sampling} draws the negative class in proportion to how strongly it
currently votes rather than uniformly, so the machine corrects the classes it actually
confuses first --- worth having when $K$ is large.

\subsection{Batched feedback: the core reformulation}
\label{sec:batched}

The classical algorithm applies feedback one example at a time, re-evaluating the clauses in
between. Each step draws random numbers for the whole $(C, 2F)$ state and writes it back. For
a machine with 5\,000 clauses over 784 features that is a 7.8-million-element stochastic pass
\emph{per training example} --- on a GPU, a sequence of small kernels with nothing to
parallelise across.

\ttlibname reformulates the update around a simple observation: \emph{every stochastic
decision in Type~I feedback is an independent Bernoulli trial with a state-independent
probability}. Memorisation happens with probability $p_{\text{mem}} = (s-1)/s$ (or $1$ with
boosted true positives), forgetting with probability $1/s$. Whether a particular
(clause, literal) automaton is offered such a trial depends only on which feedback events the
batch produced --- not on the automaton's own state. So the trials can be \emph{counted}
first and \emph{sampled} once.

Let $S_{\mathrm{I}}, S_{\mathrm{II}} \in \{0,1\}^{B\times C}$ mark which (example, clause)
pairs receive Type~I and Type~II feedback among the clauses that \emph{matched}, and
$S_{\mathrm{I}}^{\text{nofire}}$ those that received Type~I without matching. Then, with
$L \in \{0,1\}^{B\times 2F}$ the batch literal matrix:

\begin{align}
  n^{\text{true}} &= S_{\mathrm{I}}^{\!\top} L, &
  n^{\text{false}} &= S_{\mathrm{I}}^{\!\top} (1 - L), \\
  n^{\text{Ib}}_j &= \textstyle\sum_b S^{\text{nofire}}_{\mathrm{I}}[b,j], &
  n^{\text{II}} &= S_{\mathrm{II}}^{\!\top} (1 - L).
\end{align}

Four count tensors, obtained from two matrix products and a column sum. The transition then
follows in closed form:

\begin{equation}
  \Delta \;=\; \underbrace{\mathrm{Bin}\!\left(n^{\text{true}},\, p_{\text{mem}}\right)}_{\text{Type Ia memorise}}
  \;-\; \underbrace{\mathrm{Bin}\!\left(n^{\text{false}} + n^{\text{Ib}},\, \tfrac{1}{s}\right)}_{\text{Type Ia forget } + \text{ Type Ib}}
  \;+\; \underbrace{\min\!\left(n^{\text{II}},\, N - \theta\right)}_{\text{Type II, bounded at the boundary}},
  \label{eq:delta}
\end{equation}

and $\theta \leftarrow \mathrm{clip}(\theta + \Delta,\, 0,\, 2N-1)$. The two forgetting
sources merge because both are Bernoulli$(1/s)$ trials and a sum of independent Bernoulli
trials with equal probability is binomial. Type~II is deterministic, and the $\min$ reproduces
a property of the sequential algorithm exactly: an excluded False literal is pushed towards
inclusion one step at a time, but once it crosses the boundary the clause no longer matches
and receives no further Type~II events in that update.

\begin{figure}[htbp]
\centering
\resizebox{\linewidth}{!}{% Figure: sequential per-example commits vs the batched event-counting formulation.
\begin{tikzpicture}[font=\sffamily\small]

% ===================================================================== A: sequential
\node[anchor=north west, font=\sffamily\bfseries\small, text=ttslate] at (0,0.75)
  {A\ \ Classical sequential algorithm \textemdash\ $B$ stochastic commits per mini-batch};

\foreach \i/\lbl in {0/{$x_1$},1/{$x_2$},2/{$x_3$}}{
  \pgfmathsetmacro{\x}{\i*4.45}
  \node[ttnode slate, anchor=north west, text width=3.7cm, align=center] at (\x,0)
    {\lbl\ \textemdash\ evaluate $\to$ select\\[2pt]
     \textbf{draw} $(C,2F)$ randoms\\[1pt]
     \textbf{write} the automata};
}
\node[anchor=north west, font=\sffamily, text=ttgray] at (13.5,-0.55) {$\cdots$};

\draw[ttarrow, draw=ttslate] (4.05,-0.65) -- (4.5,-0.65);
\draw[ttarrow, draw=ttslate] (8.5,-0.65) -- (8.95,-0.65);
\draw[ttarrow, draw=ttslate] (12.95,-0.65) -- (13.4,-0.65);

\node[ttlabel, anchor=north west, text width=14.3cm, font=\sffamily\small] at (0,-1.75)
  {Each example sees the automata the previous example left behind. That serial dependency is
   the \emph{point} of the algorithm --- and the reason a GPU sits idle: the work between two
   state writes is far too small to fill it.};

% ===================================================================== B: batched
\node[anchor=north west, font=\sffamily\bfseries\small, text=ttdeep] at (0,-3.35)
  {B\ \ Batched formulation \textemdash\ count the events, then commit \emph{once}};

\node[ttnode, anchor=north west, text width=1.85cm, align=center] (xb) at (0,-4.05)
  {$x_{1..B}$\\[1pt]\scriptsize $(B,F)$};
\node[ttnode, anchor=north west, text width=2.25cm, align=center] (lit) at (2.45,-4.05)
  {literals\\[1pt]\scriptsize $(B,2F)$};
\node[ttnode, anchor=north west, text width=2.65cm, align=center] (cl) at (5.35,-4.05)
  {clause outputs\\[1pt]\scriptsize $(B,C)$, one matmul};
\node[ttnode, anchor=north west, text width=2.45cm, align=center] (vt) at (8.65,-4.05)
  {vote sums\\[1pt]\scriptsize $(B,K)$, clipped};
\node[ttnode, fill=ttamber!25, draw=ttamber, anchor=north west, text width=2.85cm, align=center]
  (sl) at (11.75,-4.05) {feedback masks\\[1pt]\scriptsize $S_{\mathrm{I}}, S_{\mathrm{II}} \in \{0,1\}^{B\times C}$};

\draw[ttarrow] (2.0,-4.6) -- (2.4,-4.6);
\draw[ttarrow] (4.9,-4.6) -- (5.3,-4.6);
\draw[ttarrow] (8.2,-4.6) -- (8.6,-4.6);
\draw[ttarrow] (11.3,-4.6) -- (11.7,-4.6);

% ---- counting
\node[ttnode, anchor=north west, text width=13.9cm, align=center] (cnt) at (0.35,-5.95)
  {\textbf{event counts} \textemdash\ two matrix products over the chunk\\[4pt]
   $n^{\text{true}} = S_{\mathrm{I}}^{\!\top} L$ \qquad
   $n^{\text{false}} = S_{\mathrm{I}}^{\!\top}(1-L)$ \qquad
   $n^{\text{Ib}}_j = \textstyle\sum_b S^{\text{nofire}}_{\mathrm{I}}[b,j]$ \qquad
   $n^{\text{II}} = S_{\mathrm{II}}^{\!\top}(1-L)$};
\draw[ttarrow] (13.2,-5.25) .. controls (14.2,-5.4) and (14.2,-5.6) .. (13.2,-5.9);

% ---- commit
\node[ttnode dark, anchor=north west, text width=13.9cm, align=center] (com) at (0.35,-7.85)
  {\textbf{one commit} \textemdash\ the counts become automaton transitions\\[5pt]
   $\Delta \;=\; \mathrm{Bin}\!\left(n^{\text{true}},\, p_{\text{mem}}\right)
    \;-\; \mathrm{Bin}\!\left(n^{\text{false}} + n^{\text{Ib}},\, \tfrac{1}{s}\right)
    \;+\; \min\!\left(n^{\text{II}},\, N - \theta\right)$\\[5pt]
   $\theta \;\leftarrow\; \mathrm{clip}\!\left(\theta + \Delta,\; 0,\; 2N-1\right)$};
\draw[ttarrow, line width=1.1pt] (7.3,-7.25) -- (7.3,-7.8);

\node[ttlabel, anchor=north west, text width=14.3cm, font=\sffamily\small] at (0,-10.35)
  {Every stochastic decision of the sequential algorithm is an independent Bernoulli trial, so
   the trials one automaton receives in a batch sum to a \emph{binomial} draw --- one
   $(C,2F)$ random tensor per batch instead of one per example. At $B=1$ this reduces exactly
   to the classical algorithm.};
\end{tikzpicture}
}
\caption{The reformulation. The sequential algorithm (A) performs one stochastic $(C,2F)$
  pass per example. The batched algorithm (B) accumulates the same events as counts and
  realises them in a single binomial draw per batch. Cost per update becomes almost
  independent of $B$, which is exactly the property a GPU rewards.}
\label{fig:batched}
\end{figure}

\paragraph{What is exact and what is not.}
At $B = 1$ Equation~\eqref{eq:delta} reduces to the classical algorithm term by term. For
$B > 1$ it is an approximation in precisely the sense mini-batch SGD approximates SGD: the
examples within a batch do not see each other's updates. Two secondary effects follow ---
an automaton can move by more than one step in a commit, and the Type~II bound is evaluated
against the state at commit time rather than after each event. Section~\ref{sec:fidelity}
measures what this costs. For work where fidelity matters more than throughput,
\code{feedback\_mode="sequential"} (or \code{update(..., sequential=True)}) runs the exact
per-example algorithm, still vectorised across clauses and literals.

\begin{ttcallout}[A design decision that measurement reversed]
An obvious refinement is to \emph{thin} the accumulated events --- to cap how many Type~Ia or
Type~II events one clause may contribute per batch, on the grounds that the sequential
algorithm would have re-evaluated in between and delivered fewer. This was implemented and
measured; it made learning distinctly worse. Aggregating all events is therefore deliberate,
not an oversight, and is documented as such in the repository so that the next person to have
the idea finds the answer before spending a week on it.
\end{ttcallout}

\subsection{Memory as an explicit budget}

Intermediates, not the model, dominate memory. \ttlibname never silently allocates a tensor
proportional to $B$ without a bound: each model declares how many elements one example costs,

\[
  \text{flat: } 4\,(C + 2F),
  \qquad
  \text{convolutional: } 3PC + 2C\cdot 2F + P\cdot 2F ,
\]

where $P$ is the number of patches per image, and the base class splits the batch into chunks
that keep the largest intermediate under \code{max\_chunk\_elements} (default $2^{27}$
elements). Chunking never changes the result of a batched update --- accumulation is pure
addition --- so it is purely a memory-for-nothing trade. Getting the per-example estimate
wrong is not benign, however: an under-estimate produces many tiny chunks per batch and a
4--20$\times$ slowdown, which is why the estimate is part of the subclass contract rather
than a heuristic.

\subsection{Convolutional machines}
\label{sec:conv}

A convolutional Tsetlin machine~\ttcite{2} lets a clause look for a pattern \emph{anywhere} in
an image instead of memorising absolute positions. Each $k_h\times k_w$ window becomes a patch
whose Boolean features are the $Z k_h k_w$ pixels followed --- when
\code{position\_encoding=True} --- by thermometer-encoded window coordinates, $P_y-1$ bits
$[y > k]$ and $P_x-1$ bits $[x > k]$. A clause is True for the image if it matches at least
one patch, and a feedback event updates the clause from \emph{one uniformly random matching
patch}.

\begin{figure}[htbp]
\centering
\resizebox{\linewidth}{!}{% Figure: convolutional evaluation — patches, OR over patches, and where the random patch is drawn.
\begin{tikzpicture}[font=\sffamily\small]

  \def\g{0.32}   % pixel size

  % ---------------------------------------------------------------- input image
  \node[tttitle, anchor=south west, font=\sffamily\bfseries\small] at (0,0.55)
    {Boolean image};
  \foreach \r in {0,...,5}{
    \foreach \c in {0,...,5}{
      \pgfmathsetmacro{\on}{mod(\r*3+\c*2,5)<2 ? 1 : 0}
      \ifdim \on pt > 0.5pt
        \fill[ttorange!75] (\c*\g,-\r*\g) rectangle ++(\g,\g);
      \fi
      \draw[ttgray!35, line width=0.3pt] (\c*\g,-\r*\g) rectangle ++(\g,\g);
    }
  }
  \draw[ttdeep, line width=1.3pt] (1*\g,-1*\g) rectangle ++(3*\g,3*\g);
  \draw[ttdeep!45, line width=0.9pt, dashed] (2*\g,-2*\g) rectangle ++(3*\g,3*\g);

  % ---------------------------------------------------------------- patch matrix
  \node[tttitle, anchor=south west, font=\sffamily\bfseries\small] at (3.0,0.55)
    {patch literals $(B,P,2F)$};
  \begin{scope}[shift={(3.0,0)}]
    \foreach \r in {0,...,4}{
      \foreach \c in {0,...,9}{
        \pgfmathsetmacro{\on}{mod(\r*5+\c*3,7)<3 ? 1 : 0}
        \ifdim \on pt > 0.5pt \fill[ttorange!60] (\c*\g,-\r*\g) rectangle ++(\g,\g); \fi
        \draw[ttgray!30, line width=0.3pt] (\c*\g,-\r*\g) rectangle ++(\g,\g);
      }
    }
  \end{scope}
  \draw[ttarrow] (2.15,-0.8*\g) -- (2.85,-0.8*\g);
  \node[ttlabel, anchor=south, font=\sffamily\scriptsize] at (2.5,-0.75*\g) {\code{unfold}};

  % ---------------------------------------------------------------- match tensor
  \node[tttitle, anchor=south west, font=\sffamily\bfseries\small] at (7.2,0.55)
    {matches $(B,P,C)$};
  \begin{scope}[shift={(7.2,0)}]
    \foreach \r in {0,...,4}{
      \foreach \c in {0,...,5}{
        \pgfmathsetmacro{\on}{mod(\r*4+\c*5,9)<2 ? 1 : 0}
        \ifdim \on pt > 0.5pt
          \fill[ttdeep] (\c*\g,-\r*\g) rectangle ++(\g,\g);
        \else
          \fill[white] (\c*\g,-\r*\g) rectangle ++(\g,\g);
        \fi
        \draw[ttgray!30, line width=0.3pt] (\c*\g,-\r*\g) rectangle ++(\g,\g);
      }
    }
  \end{scope}
  \draw[ttarrow] (6.35,-0.8*\g) -- (7.05,-0.8*\g);

  % ---------------------------------------------------------------- OR over patches
  \node[ttnode dark, anchor=west, text width=2.5cm, align=center] (orr) at (9.6,-0.4)
    {$c_j(X) = \bigvee_{p} m_{p,j}$\\[2pt]\scriptsize the clause fires if \emph{any} patch matches};
  \draw[ttarrow] (9.2,-0.4) -- (9.55,-0.4);

  \node[ttnode, anchor=west, text width=1.85cm, align=center] (vot) at (12.75,-0.4)
    {votes $(B,K)$\\[1pt]\scriptsize unchanged};
  \draw[ttarrow] (12.35,-0.4) -- (12.7,-0.4);

  % ---------------------------------------------------------------- sub-labels
  \node[ttlabel, anchor=north west, text width=2.6cm, font=\sffamily\scriptsize] at (0,-2.05)
    {$(Z,H,W)$ planes; a $k_h\times k_w$ window with stride $s$ gives $P$ positions};
  \node[ttlabel, anchor=north west, text width=3.2cm, font=\sffamily\scriptsize] at (3.0,-2.05)
    {$Z k_h k_w$ pixel bits, then thermometer position bits $[y>k]$, $[x>k]$};
  \node[ttlabel, anchor=north west, text width=2.1cm, font=\sffamily\scriptsize] at (7.2,-2.05)
    {one matmul over all (patch, clause) pairs};

  % ---------------------------------------------------------------- feedback path
  \node[ttnode, fill=ttamber!22, draw=ttamber, anchor=north west, text width=14.0cm]
    (fb) at (0,-3.75)
    {\textbf{Feedback} \textemdash\ for each selected (example, clause) event, \emph{one
     uniformly random matching patch} is drawn from the match tensor and its literals are
     counted: \code{p = argmax(rand(n,P) * cand)} over the candidate mask \code{cand}, then
     \code{index\_add\_} into the accumulators. Type~Ib needs no patch at all.};

  \draw[ttarrow dashed, draw=ttamber!85!black, line width=1.0pt] (11.0,-1.25) -- (11.0,-3.7);
  \node[ttlabel, anchor=west, text=ttdeep, text width=3.2cm, font=\sffamily\scriptsize]
    at (11.15,-2.4) {the match tensor is reused here, not recomputed};

  \node[ttlabel, anchor=north west, text width=14.0cm, font=\sffamily\small] at (0,-6.25)
    {Drawing the patch here rather than during evaluation keeps \emph{prediction} free of
     randomness --- a $(B,P,C)$ random tensor on every forward pass cost roughly $5\times$ the
     runtime --- and lets Type~Ia and Type~II draw independently, which matters for coalesced
     and multi-label models where one clause can receive both in the same update.};
\end{tikzpicture}
}
\caption{Convolutional evaluation. \code{unfold} produces the patch literals, one matmul
  produces all (patch, clause) matches, and an OR over patches gives the clause output. The
  random matching patch is drawn during \emph{feedback counting}, over selected events only.}
\label{fig:conv}
\end{figure}

Implementation-wise this is \code{torch.nn.functional.unfold}, one matmul, and an
\code{any} reduction. The subtlety is \emph{where} the random patch is drawn. Drawing it
during evaluation is the obvious choice and it is wrong twice over: it costs a $(B,P,C)$
random tensor and an \code{argmax} on \emph{every} forward pass, including pure prediction
where the draw is discarded --- roughly $5\times$ the runtime on MNIST-sized inputs --- and it
forces Type~Ia and Type~II to share one patch, which is incorrect for coalesced and
multi-label models where a clause can receive both in the same update. \ttlibname therefore
returns the per-patch match tensor as an opaque \emph{feedback context} from evaluation and
draws inside the feedback counter, over the selected events only. Prediction is RNG-free, and
a test pins that it stays so.

The random draw itself is a small trick worth noting: for candidate mask
$\mathrm{cand} \in \{0,1\}^{n\times P}$, \code{argmax(rand(n,P) * cand)} selects a uniformly
random index among the ones, with no rejection loop and no gather over a variable-length list.

\begin{ttcallout}[Position encoding is a modelling choice, not a default to leave alone]
Position literals let a clause demand \emph{where} it matches. On genuinely
translation-invariant tasks they hurt --- on the generated \code{make\_shapes} benchmark,
accuracy oscillates between 0.5 and 0.9 with them enabled and reaches 1.000 on every seed
with them disabled. On MNIST, where digits are centred and position is informative, they are
worth about a point of accuracy. Check this flag before suspecting the feedback code.
\end{ttcallout}

\subsection{Booleanization}

The literals are the vocabulary of every rule the machine can learn, so Booleanization is the
most consequential modelling decision in a Tsetlin machine project --- and in practice it
moves accuracy more than the model hyper-parameters do. \ttlibname treats encoders as
first-class modules (Table~\ref{tab:encoders}): they are \code{nn.Module}s with a
\code{fit}\slash\code{transform} protocol, they move to a device, they are saved in a
\code{state\_dict}, and they compose.

\begin{table}[htbp]
\centering
\small\sffamily\sloppy
\caption{Booleanization encoders. $F$ input columns, $b$ bits per column.}
\label{tab:encoders}
\begin{tabularx}{\linewidth}{@{}L{4.4cm}L{2.15cm}>{\raggedright\arraybackslash}X@{}}
\ttoprule
\thead{Encoder} & \thead{Output} & \thead{Use and why}\\
\ttmidrule
\clsname{Binarizer} & $F$ &
  a fixed threshold per element, the original MNIST recipe (\code{pixel > 0.3})\\
\clsname{ThermometerEncoder} & $F\!\cdot\!b$ &
  ordinal bits $[v \ge t_1,\dots,v \ge t_b]$; two literals then express an interval
  (\code{age>=30 AND NOT age>=50}). Thresholds by quantile, uniform width or unique value\\
\clsname{OneHotEncoder} & $\sum_i n_i$ & one bit per category; rules read as \code{colour=red}\\
\clsname{BitPlaneEncoder} & $F\!\cdot\!b$ & place-value bits of integers; compact codes\\
\clsname{AdaptiveThresholdEncoder} & $F$ &
  local Gaussian\slash mean thresholding, the standard recipe for MNIST-like images in the
  convolutional TM literature\\
\clsname{ColorThermometerEncoder} & $C\!\cdot\!b$ planes & per-channel thermometer for RGB\\
\clsname{HypervectorEncoder} & $d$ &
  sparse random hypervectors bundled with OR, optionally binding position by cyclic shift ---
  a compact encoding for large vocabularies~\ttcite{13}\\
\ttbotrule
\end{tabularx}
\end{table}

Encoders also carry the inverse mapping, which is what keeps rules readable. A thermometer
encoder's \code{decode\_range} turns a clause's included bits back into an interval in the
original units, and \code{aggregate\_literal\_importance} sums literal importances back over
the bits of one original column.

\begin{ttcallout}[A GPU bug worth documenting]
Encoders keep their thresholds on the CPU, so handing one a CUDA tensor triggers a
device-to-host copy. Issued with \code{non\_blocking=True} into pageable memory, such a copy
returns \emph{before} the data lands, and the encoder thresholds stale memory --- about 40\%
wrong bits, no error, just a quietly worse model. All coercion in \ttlibname routes through a
helper that sets \code{non\_blocking} only for CUDA destinations. The regression test has to
allocate on the device (\code{torch.rand(..., device=dev)}) rather than moving a host tensor
across, or the race never occurs and the test passes vacuously.
\end{ttcallout}

\subsection{Interpretation}

Three levels of explanation are available, all closed-form and all computed from the state
rather than by probing the model.

\paragraph{Rules.} \code{model.rules()} prints each non-empty clause as
\code{IF <conjunction> THEN <class>}, with its weight where weights exist. Feature names flow
through from the encoder, so the conjunction is in the user's vocabulary.

\paragraph{Global importance.} Following~\ttcite{4}, the strength of literal $l$ for output
$k$ is the weighted fraction of positive clauses of $k$ that include it,
\[
  g_k[l] \;=\; \frac{1}{m_k}\sum_{j \,:\, W_{jk} > 0} W_{jk}\,\mathrm{include}_{jl},
  \qquad m_k = \left|\{\,j : W_{jk} > 0\,\}\right| .
\]

\paragraph{Local importance and explanations.} For one example, the same quantity restricted
to the clauses that actually matched and to the literals that are True gives a per-example
attribution; \code{interpret.explain} returns the matching clauses ranked by weight, which is
an account of the prediction in the model's own terms rather than a surrogate.

\subsection{Correctness}

The library is tested with 44 tests, parametrised over a \code{device} fixture so that every
device-marked test runs twice on a GPU machine. The tests that matter most are not the
smoke tests but the semantic pins: brute-force reference implementations of clause evaluation
and of convolutional patch semantics; the empty-clause train\slash eval asymmetry; the
RNG-free prediction path for convolutional models; the encoder device-transfer race described
above; and a packaging test asserting that the version agrees between
\code{pyproject.toml} and \code{\_\_init\_\_.py} and that every subpackage ships in the wheel
--- which exists because a \code{.gitignore} rule once silently excluded
\code{torchtsetlin/data} from a release.

\begin{ttkey}
Clause evaluation is a matmul against the include mask; learning is event counting plus one
binomial draw per commit. Those two identities are the whole library: everything else ---
models, convolution, chunking, the trainer --- is organisation around them.
\end{ttkey}
